% from ega3-6-fr.tex
\oldpage[III]{148}
déduit par fonctorialité des homomorphismes canoniques
$P_\bullet\to(P_\bullet)_f$, $Q_\bullet\to(Q_\bullet)_g$, est bijectif. Or
($0_{\mathrm I}$, 1.6.1), on peut écrire
$(P_\bullet)_f=\varinjlim P_\bullet^{(n)}$, où les $P_\bullet^{(n)}$ sont tous
des complexes de $B$-modules identiques à $P_\bullet$, l'application
$P_\bullet^{(m)}\to P_\bullet^{(n)}$ pour $m\leqslant n$ étant la
multiplication par $f^{n-m}$; on a un résultat analogue pour $Q_\bullet$ en
remplaçant $f$ par $g$; d'autre part, il est clair que l'homomorphisme
\[
  \mathbf{Tor}_\bullet^A(P_\bullet^{(m)},Q_\bullet^{(m)})
  \longrightarrow
  \mathbf{Tor}_\bullet^A(P_\bullet^{(n)},Q_\bullet^{(n)})
\]
correspondant aux homomorphismes
$P_\bullet^{(m)}\to P_\bullet^{(n)}$ et
$Q_\bullet^{(m)}\to Q_\bullet^{(n)}$ est par définition la multiplication
par $(f\otimes g)^{n-m}$. La conclusion résulte alors de ($0_{\mathrm I}$,
1.6.1) appliqué au premier membre de (6.4.5.3) et de (6.3.6).
\end{proof}

\begin{env}[6.4.6]
Avec les notations de (6.4.4), on définit de même des homomorphismes canoniques
de foncteurs spectraux
\[
  \left\{
  \begin{aligned}
  (u\times_S v)^*\bigl('\mathscr E(\mathscr P_\bullet,\mathscr Q_\bullet)\bigr)
    &\longrightarrow
      '\mathscr E(u^*(\mathscr P_\bullet),v^*(\mathscr Q_\bullet)),\\
  (u\times_S v)^*\bigl(''\mathscr E(\mathscr P_\bullet,\mathscr Q_\bullet)\bigr)
    &\longrightarrow
      ''\mathscr E(u^*(\mathscr P_\bullet),v^*(\mathscr Q_\bullet)).
  \end{aligned}
  \right.
  \tag{6.4.6.1}\label{III.6.4.6.1.fr}
\]
et le raisonnement de (6.4.5) montre que lorsque $u$ et $v$ sont des
immersions ouvertes, les homomorphismes (6.4.6.1) sont bijectifs : en effet,
compte tenu de (6.3.6.2) et (6.3.6.3), il prouve que c'est un isomorphisme pour
les termes $E^2$, et (6.4.5) montre que c'est un isomorphisme pour les
aboutissements; on conclut donc à l'aide de (0, 11.1.2) et (0, 11.2.4).
\end{env}

\subsection{Foncteurs hypertor locaux de complexes de Modules quasi-cohérents : cas général}

\begin{env}[6.5.1]
Considérons maintenant un préschéma $S$ quelconque et deux $S$-préschémas
quelconques $X$, $Y$; soit $\mathscr P_\bullet$ (resp. $\mathscr Q_\bullet$)
un complexe de $\mathscr O_X$-Modules (resp. $\mathscr O_Y$-Modules)
quasi-cohérents. Posons $Z=X\times_S Y$; nous allons définir des
$\mathscr O_Z$-Modules quasi-cohérents
$\mathscr{T}\!or_n^S(\mathscr P_\bullet,\mathscr Q_\bullet)$ dits
\emph{hypertor locaux} de $\mathscr P_\bullet$ et $\mathscr Q_\bullet$ qui se
réduiront à ceux déjà définis dans (6.4) lorsque $S$, $X$ et $Y$ sont affines.

Lorsque $\mathscr P_\bullet$ et $\mathscr Q_\bullet$ se réduisent
respectivement à leurs termes de degré 0, $\mathscr F$ et $\mathscr G$ (les
autres étant nuls), on écrira
$\mathscr{T}\!or_n^S(\mathscr F,\mathscr G)$ au lieu de
$\mathscr{T}\!or_n^S(\mathscr P_\bullet,\mathscr Q_\bullet)$.
\end{env}

\begin{env}[6.5.2]
Supposons d'abord $S$ affine, et soient $(X_\lambda)$, $(Y_\mu)$ des
recouvrements de $X$ et $Y$ respectivement par des ouverts affines; alors les
$Z_{\lambda\mu}=X_\lambda\times_S Y_\mu$ forment un recouvrement ouvert affine
de $Z$. Posons
$\mathscr P_{\lambda\bullet}=\mathscr P_\bullet|X_\lambda$,
$\mathscr Q_{\mu\bullet}=\mathscr Q_\bullet|Y_\mu$; nous avons donc pour tout
couple $(\lambda,\mu)$ un $\mathscr O_{Z_{\lambda\mu}}$-Module quasi-cohérent
\[
  \mathscr F_{\lambda\mu}
  =\mathscr{T}\!or_n^S
   (\mathscr P_{\lambda\bullet},\mathscr Q_{\mu\bullet}),
\]
et il faut montrer que les $\mathscr F_{\lambda\mu}$ vérifient la condition de
recollement ($0_{\mathrm I}$, 3.3.1). Pour cela, il suffit de vérifier que pour
tout ouvert affine $U\subset X_\lambda\cap X_{\lambda'}$ (resp.
$V\subset Y_\mu\cap Y_{\mu'}$), les restrictions de
$\mathscr F_{\lambda\mu}$ et de $\mathscr F_{\lambda'\mu'}$ à $U\times_S V$
sont canoniquement isomorphes; mais cela découle aussitôt de l'existence
d'isomorphismes canoniques de ces restrictions sur
$\mathscr{T}\!or_n^S(\mathscr P_\bullet|U,\mathscr Q_\bullet|V)$ (6.4.5). En
outre, il résulte aussitôt de cette définition et de (6.4.5) que le
$\mathscr O_Z$-Module ainsi défini ne dépend pas (à un isomorphisme près) des
recouvrements ouverts consi